\section{Descent from the marked field to the two ordinary levels}
\label{sec:host-descent}

The preceding appendix produced integral points on the larger marked field. Which
constituents can they supply at the smaller ordinary levels where we will evaluate
Frobenius traces? We determine the irreducible constituents
after restriction and keep track of the ordinary flags and
coefficient types in descent. The two possibilities for a
rank-three point give different packets in the length
calculation that follows. Lemma~\ref{lem:rank-three-host-restriction} controls which constituents remain irreducible on restriction. Proposition~\ref{prop:ordinary-host-attachment} returns the corresponding cuspidal systems to their original fields.

Retain the quadratic field \(K\), the marked element and the support conditions of Section~\ref{sec:marked-hosts}. Write \(H=G_K\) for the kernel of the quadratic character defining \(K/\mathbb Q\). At a type-I point the rank-three representation is irreducible on \(H\). At a type-II point for the general datum, let \(W\) be the outer plane and \(U\) the middle line. On \(H\) the preserved decomposition is
\[
 V=W\oplus U,\qquad \dim W=2,\quad\dim U=1,\quad W|_H\text{ irreducible}.
\] The length argument distinguishes the packets
coming from these two cases. We first determine the group on
which irreducibility actually holds.

\paragraph{Restriction of the rank-three points.}
Start with restriction, where an irreducible representation may acquire several
constituents. Use the integral representative of a type-I point and its
ordinary flag at \(B\). The residual trace on \(H\) is
\(P+1+P^{-1}\). Given an open subgroup \(V\) of \(H\), its normal core is the kernel of the finite permutation action:
\[
 U=\bigcap_{h\in H}hVh^{-1}=\ker\bigl(H\longrightarrow\operatorname{Perm}(H/V)\bigr).
\] In \(U\cap I_B\) choose \(u\)
with \(\psi(u)=1\) and \(\epsilon(u)\) of infinite order.
Indeed rational cyclotomic inertia surjects onto
\(\mathbb Z_\ell^\times\). Start with a lift of \(1+\ell\)
and take a power to enter the kernels of \(\eta,\psi\) and
the open core. In an ordinary basis the three eigenvalues
of \(u\) are \(\epsilon(u)^e,1,\epsilon(u)^{-e}\). A further
upper triangular change makes this matrix diagonal.

\begin{lemma}[Absolute irreducibility after restriction]
\label{lem:rank-three-host-restriction}
Let \(r\) be a type-I rank-three point with the integral
ordinary flag at \(B\) and residual trace \(P+1+P^{-1}\) on the
quadratic kernel \(H\) of Appendix~\ref{sec:linked-seeds}.
Assume the marked cyclotomic and local inertia data at that level.
The two restriction statements are
\[
 U\le H\text{ open}\ \Longrightarrow\ r|_U\text{ absolutely irreducible},
 \qquad r|_{\ker(\epsilon|_H)}\text{ absolutely irreducible}.
\] The associated
rank-two corner at the two-block points is treated separately
in the ordinary support comparison.
\end{lemma}

\begin{proof}
\proofstep{The isotypic summands after restriction.}
If restriction to \(U\) is reducible, it is semisimple because
the sum of the \(H\)-translates of a simple \(U\)-subspace
is the whole space, and a sum of simples is semisimple.
The isotypic summands form a single orbit. Since three
is prime, either \(U\) acts by scalars or there are three
distinct character lines permuted transitively. The first
case contradicts the spectrum of \(u\). In the second case
the lines are precisely the three eigenlines of \(u\).
Every matrix of \(H\) is monomial in this basis. The
separating element \(x\in D_B\cap H\) constructed above is
also upper triangular, so fixes each line. A transitive
subgroup of \(S_3\) contains a three-cycle.

For a lift \(g\)
of this cycle, \(gx^n\) has trace zero for every \(n\).
Reducing the trace in the original integral model and writing
\(a=P(g)\), \(b=P(x)\) gives
\[
         (ab^n)^2+ab^n+1=0\qquad(n=0,1,2).
\]
The three roots are distinct since \(b^2\ne1\), a
contradiction in every residue characteristic. Thus the
representation \(\rho\) is irreducible on every open
subgroup of \(H\), in particular on \(G_L\).

\proofstep{Restriction to the full cyclotomic kernel.}
It remains irreducible on \(H\cap\ker\epsilon\). Here is a
characteristic-zero argument for this further restriction.
Let \(G\) be the Zariski closure of \(\rho(H)\) over the
algebraically closed coefficient field. The preimages of
its finitely many components are closed and open, so its
identity component \(G^\circ\) acts irreducibly. Let \(U\) be the unipotent radical and \(V\) the representation space. Unipotence gives a nonzero fixed vector and normality gives stability, so irreducibility implies
\[
 0\ne V^U\subseteq V,\qquad G^\circ V^U=V^U\quad\Longrightarrow\quad V^U=V .
\] Hence the radical is trivial.
Schur's lemma makes the center scalar. As a connected
reductive group is the product of its connected center and
derived group, the derived group acts irreducibly. The
Zariski closure of \(\rho(H\cap\ker\epsilon)\) contains it since
commutators of the dense subgroup belong to that closure
and the commutator map is algebraic. This proves the claim.

\end{proof}

\begin{proposition}[Attachment on the original fields]
\label{prop:ordinary-host-attachment}
For the contributing type-I and outer-plane points, the
ordinary support comparison and cyclic descent yield the
cuspidal systems on the designated original fields, with the
specified shifts of weight and local level. In rank three the weight shift and multiplier are
\[
 r_3=\rho\otimes\epsilon^{-1},\qquad
 r_3(cgc)^{\mathsf T}Jr_3(g)=\epsilon(g)^{-2}J\quad(g\in G_K),
\]
where \(J\) reverses the marked basis. Each cyclic step uses irreducibility on the restricted group and the stipulated
local ordinary condition.
\end{proposition}

\begin{proof}
\proofstep{The ordinary cusps over \(L\).}
At an integral point, pass to a common finite coefficient
DVR with the seed before applying the support comparisons.
On \(G_{L^+}\) the point kills the required tame kernels
and has the prescribed flag characters on inertia, so it
defines a point of the framed ordinary quotient. In rank
three the restrictions of the positive-component matrices
agree with \(\rho|_{G_L}\). In rank two they agree with the
outer block at a type-II point. Choose a second omitted-prime pair whose omitted prime differs in each rank. Writing \(p_n,p_n'\) for the first and second choices in rank \(n\), the requirement is
\[
 p_2\ne p'_2,\qquad p_3\ne p'_3 .
\]
The comparison for the first pair gives the finite fiber
used in \eqref{eq:tag-h3}. Each run also puts the irreducible point of
the relevant rank on its own Hecke support.

A common
algebraic character twist makes its inertia weights a
classical specialization of the ordinary module. A
contributing ordinary eigenform then has attached
characteristic polynomials equal to those of the point
with that twist. A decomposition into proper cuspidal or Speh blocks would split the semisimple Galois representation at matching Frobenii into pieces of smaller rank. In particular its characteristic polynomial there would be the product
\[
 \det(1-X\rho^{\rm ss}(g))=\prod_i\det(1-X\rho_i(g)),\qquad \sum_i\dim\rho_i=\dim\rho,
\]
for the global block constituents \(\rho_i\). Irreducibility excludes this. Thus the
transfer is cuspidal. Equality at the common spherical
places and strong multiplicity one identify the two
transfers after adjusting the common twist. At either
omitted prime use the local comparison from the other
run. This gives the local support and size conditions of
the ordinary comparison wherever they are required.
Undoing an algebraic character twist translates the common
center of the size interval by its absolute-value exponent.

\proofstep{The weight shift of the rank-three system.}
In rank three write \(r_3=\rho\epsilon^{-1}\). On the negative
component the frame acts by inverse transpose relative to
the reversing matrix \(J\), and the matrix part at \(c\) is
one. On \(G_K\) this gives
\[
             r_3(cgc)^t J r_3(g)=\epsilon(g)^{-2}J .
\]
The twist thus gives the multiplier \(\epsilon^{1-3}\).
It preserves every comparison of ratios of graded
characters on a flag. The algebraic weights on inertia
become \((e-1,-1,-e-1)\).

\proofstep{The rank-two corner over \(\mathbb Q\).}
For general data the character on the middle line reduces to
one on \(H\). On inertia where \(P(h)\ne1\) it is a root of
the polynomial with roots \(\Psi(h),1,\Psi(h)^{-1}\). Only
the middle root reduces to one there. At a place above
\(\ell\), choose \(s\) in inertia with \(P(s)^2\ne1\).
When \(P(h)=1\), use multiplicativity on \(s\) and \(hs\). A representation matrix at \(h\) is recovered from those two by
\[
 B(h)=B(hs)B(s)^{-1} .
\]

At the other places
the inertia with \(P=\eta=1\) belongs to the prescribed tame
kernel. Thus the middle character is unramified everywhere.
It is one modulo the maximal ideal, and the class number
of the chosen \(K\) is prime to \(\ell\). The unramified
class field quotient forces the character to be one. The
determinant of the outer block has residue one, and is one
on every inertia group by the same three-root polynomial.
The same class group argument gives determinant one on \(H\).

\proofstep{Extension of the outer plane to the rational group.}
Write \(B(g)\) for this outer block and
\(D=\operatorname{diag}(-1,1)\). For a determinant-one
two by two matrix, inverse transpose conjugated by the
reversing matrix is conjugation by \(D\), as follows by
the two by two inverse formula. Thus polarization gives \(B(cgc)=D B(g)D^{-1}\).
With \(D^2=1\), the formula
\[
 W(hc^j)=B(h)D^j,\qquad
     h\in H,\ j\in\{0,1\},
\]
defines a continuous representation on \(H\rtimes\langle c\rangle\):
the product of the two matrices uses
\(D^j B(h') D^{-j}=B(c^jh'c^{-j})\).
Its determinant is one on \(H\) and \(-1\) at \(c\), hence is
\(\eta\).

It remains irreducible on \(G_L\). Otherwise the
normal restriction is semisimple and the marked element,
whose two residual eigenvalues differ, excludes scalar
restriction. There are two distinct lines, swapped by
some \(g\in H\). Then \(g\) and \(gg_*\) both have trace
zero. Reduction on the outer block forces
\[
 P(g)^2=P(gg_*)^2=-1,
\]
contradicting \eqref{eq:tag-h1}.
The rank-two run above consequently gives a regular
cuspidal system over \(L\) for this particular corner
after its algebraic twist.

\proofstep{Descent of the global system.}
Filter the abelian extension \(L/K\) by cyclic steps of
prime degree. Each intermediate field contains the
imaginary quadratic field and is abelian over \(\Q\),
hence is CM. At a step \(F'/F\) the given representation
exists on \(G_F\) and is irreducible even on \(G_{F'}\).
If \(\sigma\) generates the cyclic group, the good Frobenius polynomials of the attaching cusp \(\pi\) and of its conjugate agree. By strong multiplicity one this gives
\[
 {}^\sigma\pi\simeq\pi .
\] Cyclic
base change descends it to a cusp on \(F\). The infinity
parameter is regular algebraic because at a complex place the
local extension is trivial and the compatible parameters
agree. The Galois representation attached to this
descent agrees after restriction with the given one by
compatibility at split places, Chebotarev and
Brauer--Nesbitt.

Here degree-one primes over \(F\) are
dense for the required Frobenius comparison over \(F'\) since
the higher relative degrees have density zero.
Identify the two restrictions on \(G_{F'}\) using the
intertwiner. If the resulting representations are
\(r,\widetilde r\), then \(\widetilde r(g)^{-1}r(g)\)
commutes with their common action of \(G_{F'}\): both matrices
carry the \(h\)-action to the \(ghg^{-1}\)-action.
Schur's lemma gives a scalar \(\nu(g)\), and
\[
       r(g)=\nu(g)\widetilde r(g),\qquad
       \nu(gh)=\nu(g)\nu(h),\qquad \nu|_{G_{F'}}=1 .
\]
This character of the cyclic quotient gives the finite
Hecke twist that selects the specified descent.
The conjugate duality required of the descended cusp
follows by the same comparison at good places from the
given pairing. Iteration gives the rank-three polarized cusp
over \(K\) and the rank-two cusp for \(W|_{G_K}\). The
alternating rank-two pairing has multiplier \(\det W\), and
composing with \(W(c)^{-1}\) supplies the conjugate pairing.
\hfill\mbox{\cite{arthur-clozel,clozel-rajan,harris-lan-taylor-thorne}}

\proofstep{The additional rank-two descent.}
There is one further cyclic descent for the rank-two
system, from \(K\) to \(\Q\). The complex parameter has
distinct weights with integral difference \(2e>0\) and
finite determinant on the Galois side. Conjugation
invariance pairs the two characters of the complex Weil
group. With positive gap, the infinity parameter of the descent is the induction from one of the distinct conjugate Weil characters \(\xi,\xi^c\):
\[
 \operatorname{Ind}_{W_{\mathbb C}}^{W_{\mathbb R}}\xi,\qquad\xi\ne\xi^c .
\]

It therefore gives the discrete series up to twist. After the algebraic normalization
\(W(-e)\) the two cyclotomic exponents are \(0,-2e\),
the geometric convention for weight \(2e+1\). Write \(r_f\) for the classical newform representation before
choosing its descent twist. Polynomial comparison at degree-one primes
gives \(r_f|_{G_K}\simeq W(-e)|_{G_K}\). Applying the scalar
intertwiner calculation to this quadratic step gives
\[
 W(-e)\simeq r_f\otimes\nu,\qquad
 \nu\in\operatorname{Hom}(\operatorname{Gal}(K/\mathbb Q),
                      \{\pm1\})=\{1,\eta\}.
\] Keeping the particular
extension is needed for both off-diagonal Selmer summands.

\proofstep{The local condition through descent.}
Let \(\kappa_i\) be the graded characters of the flag and put
\(\delta_{ij}=\kappa_i/(\epsilon\kappa_j)\). In terms of their
images on the decomposition group, the two conditions read
\begin{equation}\label{eq:host-character-image-conditions}
 \begin{aligned}
 U_v         &\iff \delta_{ij}(G_{F_v})\ne\{1\} &&\text{for all }i,j,\\
 U_v^\infty &\iff |\delta_{ij}(G_{F_v})|=\infty &&\text{for all }i,j .
 \end{aligned}
\end{equation}
Jordan--Hölder makes the tests independent of the chosen flag. For a
finite local extension with subgroup \(H\subseteq G_{F_v}\),
\[
 [\delta_{ij}(G_{F_v}):\delta_{ij}(H)]\le [G_{F_v}:H]<\infty .
\]
Thus restriction preserves the second condition in both directions.
A nontrivial finite character can become trivial on \(H\), so the
first condition of \eqref{eq:host-character-image-conditions}
needs a separate argument on the chosen extension.

\proofstep{Ordinary flags under local restriction.}
First suppose \(U_v\) holds downstairs and the ordered
inertial characters are cyclotomic powers of strictly
decreasing exponents times finite characters of level
\(N\). The ordinary representation becomes crystalline
over the finite extension killing those inertia characters
of finite order. For the induction add one graded character
\(\kappa_j\) to a previously treated subrepresentation
\(V_0\). The representation
\(V=V_0\otimes\kappa_j^{-1}\) is de Rham downstairs by
finite descent and crystalline on that extension. Its
strictly positive cyclotomic exponents give
\(\operatorname{Fil}^0D_{\rm dR}(V)=0\).

Also \(H^0(V^*(1))=0\) by \(U_v\). Local duality and the Euler
characteristic formula give
\[
 \dim H^1(F_v,V)=[F_v:\Q_\ell]\dim V+\dim H^0(F_v,V).
\]
This equals \(\dim H_f^1(F_v,V)\) by the Bloch--Kato
exponential sequence. The extension class therefore
becomes crystalline where its ends are crystalline.
This proves the induction. The dimension calculation took
place downstairs, where the hypothesis on the characters
was imposed. \hfill\mbox{\cite{bloch-kato-finite}}

A forbidden equality of characters upstairs may be moved
to one of the frame places. The positive coset acts by
conjugation, the negative one by conjugation of the dual.
Conjugating transports an equality to the corresponding decomposition group. For duality the test is symmetric in the ordered pair of indices, since
\[
 \kappa_i^{-1}=\epsilon\,\kappa_j^{-1}
      \quad\Longleftrightarrow\quad
 \kappa_j=\epsilon\,\kappa_i .
\]
Both operations preserve potential crystallinity. The decomposition group of the resulting
frame place lies on the positive component by splitting in
\eqref{eq:tag-h2}. Its characters are precisely those of the restricted
ordinary flag, by comparison of characteristic polynomials
for all group elements. The even-level conclusion of
Section~\ref{sec:linked-components} at the potentially
crystalline point excludes the equality. This proves the
required persistence. In rank two the two exponents differ
by \(2e\ge2\), so the same assertion follows directly on
every open subgroup of inertia.

\proofstep{Nonresonance at the even level.}
There is a further use of the even-level result on the
original rank-three host. At a general-data place over \(\ell\), the local point factors through the quotient for the prescribed potentially crystalline type. If its ideal is \(I_{\rm pc}\) and \(x_v\) is the point, this says
\[
 I_{\rm pc}\subseteq\ker x_v,\qquad R_v\twoheadrightarrow R_v/I_{\rm pc}\xrightarrow{\ x_v\ }\overline{\mathbb Q}_\ell .
\] Its crystalline model on
the specified open subgroup is defined over a finite
coefficient extension by evaluating the local deformation
ring there. Conjugation and duality give the model at the
conjugate place.

If an equality occurred already over
\(K\), its restriction would occur over \(L\), and the
same frame test rules it out. A different auxiliary marked
element can be used for this test by diagonalizing the marked
difference with its distinct residues and using the paired
basis to retain the normalization at \(c\). The two
representations on \(G_K\) are isomorphic. On the
auxiliary field over \(\mathbb Q(i)\) one can test directly on the element \(x\)
at \(\ell\) of Section~\ref{sec:marked-hosts}. The three eigenvalues of the separator remain distinct in the residue field:
\[
 \{P(x),1,P(x)^{-1}\}\text{ has cardinality }3,\qquad \bar\epsilon(x)=1 .
\] An equality on two different lines would imply
\[
  \kappa_i=\epsilon\kappa_j
     \quad\Longrightarrow\quad
  \bar\kappa_i(x)=\bar\epsilon(x)\bar\kappa_j(x)=\bar\kappa_j(x)
       \qquad (i\ne j),
\]
contradicting distinctness.
For the same line the cyclotomic character is nontrivial
on local inertia. Conjugating the place preserves the test.
These arguments include \(e=1\).

For clarity, the local output glued from the two ordinary
runs gives \(U_{v'}\Rightarrow U_{v'}^\infty\) over \(L\).
For an attached cusp it gives a single real \(t\) such
that its local support at \(v'\nmid\ell\), and also at
\(v'\mid\ell\) satisfying \(U_{v'}^\infty\), consists of
characters \(\xi\) with
\[
           q_{v'}^{t-1/2}<|\xi(\varpi_{v'})|
                          <q_{v'}^{t+1/2}.
\]
The interval excludes ratios \(|\cdot|^{\pm1}\). Hence
the principal series is irreducible and its parameter
has zero monodromy. Persistence and the finite-index
criterion for infinite order descend the first assertion
about \(U^\infty\).

\proofstep{Descending the valuation interval.}
To descend the interval in a cyclic step, the upstairs
support characters on units must come from below by norm.
Away from \(\ell\) the downstairs inertia polynomial is
that of a list of global finite characters. Under the stated compatibility, equality of the inertia characteristic polynomials on the full upstairs inertia identifies the semisimplification. If \(\theta_i\) are the finite characters downstairs, at the places away from \(\ell\) this reads
\[
 r_{\rm up}|_{I_{v'}}^{\rm ss}\simeq\bigoplus_i\theta_i|_{I_{v'}} .
\] Above \(\ell\), when
the interval is needed, the representation is de Rham
and the upstairs Weil parameter is toral. Exactness of
potentially semistable comparison matches its characters
with the Weil characters of the graded lines. A de Rham
character under local reciprocity is locally algebraic.
Factor the locally algebraic character as \(\chi_{\rm alg}\chi_{\rm sm}\).

Under restriction from \(v\) to \(v'\), both factors are pulled back by the local norm \(N=N_{v'/v}\), so
\[
 (\chi_{\rm alg}\chi_{\rm sm})\circ N=(\chi_{\rm alg}\circ N)(\chi_{\rm sm}\circ N).
\]
Thus at either kind of place each upstairs character
extends by norm on units and extends on the full local
multiplicative group by choosing the value of a uniformizer. It can be
chosen by taking a root, accounting for its unit factor.
Equivalently the corresponding upstairs Weil character
extends to the downstairs Weil group. Local cyclic base
change restricts the Frobenius-semisimple parameter.
Monodromy downstairs is already zero by restriction.
Each isotypic summand upstairs is preserved downstairs,
since its character extends. Dividing by an extension,
one obtains a representation of the cyclic quotient,
which splits into characters in characteristic zero.
Thus the downstairs parameter is toral too.

For a place \(v'/v\) of residue degree \(f\), restriction of the local character is norm pullback by reciprocity:
\[
 \chi|_{G_{v'}}\ \longleftrightarrow\ \chi_v\circ N_{v'/v}.
\] For an inducing
character \(\alpha\), the norm of a uniformizer is a unit
times \(\varpi_v^f\), and the smooth character on units has
absolute value one. Consequently the upstairs bound reads
\[
 q_v^{f(t-1/2)}
   <|(\alpha\circ\mathrm N)(\varpi_{v'})|
      =|\alpha(\varpi_v)|^f
   <q_v^{f(t+1/2)} .
\]
Taking positive \(f\)-th roots gives exactly the original
bounds downstairs. Split places give the identity
comparison. These arguments iterate, and also transport
the rational rank-two comparison to an imaginary quadratic
field on which \(W\) remains irreducible, using the
newform's base change. \hfill\mbox{\cite{arthur-clozel,blggt-lg}}

These descents apply to the integral torus representatives
of Proposition~\ref{prop:marked-invariant-algebra}. The
following level calculation uses the original inertia
characters, so the resulting packet counts do not carry a
factor from the larger marked field.

\proofstep{Level on the original field.}
Away from \(\ell\) the type-I inertia is conjugate to
\(\psi,1,\psi^{-1}\) on \(H\). At \(\ell\) its flag has the same finite parts. Twisting by
\(\epsilon^{-1}\) therefore changes the ordered cyclotomic exponents by
\[
             (e,0,-e)\ \longmapsto\ (e-1,-1,-e-1).
\] For \(W\) at \(p\ne\ell\) the inertia
types on \(I_p\cap H\) are \(\psi^{\pm1}\). Each character
extends to \(I_p\), so the isotypic spaces are preserved.
If the subgroup has index two, the remaining generator
has scalar square on each such space, hence diagonalizes.
The extensions differ from the specified characters by
a power of \(\eta\). This gives the level
\(M_2=N|D_K|\).

At \(\ell\) a nonzero middle coordinate of the ordinary first line would intertwine the high inertia character \(\chi\) with the trivial one. If its value is \(a\), equivariance gives
\[
 (\chi(h)-1)a=0\quad(h\in I_\ell),\qquad \chi|_{I_\ell}\ne1\quad\Longrightarrow\quad a=0 .
\]
The first line therefore projects into the outer plane. They differ on an infinite-order
cyclotomic element after killing \(\psi\). In the plane
this inertia line is unique, since the second character
has exponent \(-e\), and it is stable under decomposition
because its character on inertia is the restriction of a
global character. For the plane, the two inertia exponents and the conductor condition on their finite factors \(\theta_i\) are
\[
 (e,-e),\qquad \operatorname{cond}(\theta_i)\mid M_2\quad(i=1,2).
\]

\proofstep{The fixed-vector bounds at the descended level.}
Here is the passage from these level data to fixed vectors.
Use \(M_3=N\) in rank three and \(M_2\) in rank two.
The ordinary de Rham jumps determine the archimedean
infinitesimal character with entries \((n-1)/2+m_i\)
in the algebraic Satake convention, where \(m_i\) are
the displayed exponents. At each finite place the local
parameter is toral with zero monodromy by the size and
nonresonance comparison. Each of its characters is trivial
on the inertia fixing \(\mu_{M_n}\). At \(\ell\) the smooth
parts of the de Rham characters give precisely this
condition.

At a place \(v\mid M_n\) over \(p\), the norm on local units satisfies
\[
 \mathrm N_{F_v/\mathbb Q_p}(1+M_n\mathcal O_v)
                    \subseteq 1+M_n\mathbb Z_p ,
\]
where \(F_v\) denotes the local field of the descended group.
Reciprocity and the inertia condition imply that each inducing
character is trivial on \(1+M_n\mathcal O_v\).
For the irreducible principal series, the induced-function
model has a vector supported on \(B K(M_n)_v\) taking
value one at the identity because the inducing character and
modulus are trivial on the intersection with the compact
subgroup. At the remaining places the unramified toral
parameter gives a spherical vector. Taking their tensor
product gives the required finite-adelic fixed vector.
For the rational packet the same argument applies after
irreducible quadratic base change wherever the packet
comparison requests it. Neither discriminants nor
principal levels of the enlarged field occur in \(M_n\).

\end{proof}

\paragraph{The rational adjoint Selmer groups.}
Having identified the possible constituents, we need to control their rational adjoint
deformations. For a regular algebraic polarized cusp over a CM field
\(F\), Thorne's adjoint vanishing theorem applies if the
attached representation is irreducible on
\(G_{F(\mu_{\ell^\infty})}\). It gives vanishing of
\(H_f^1(F^+,\operatorname{ad}r)\) for the extension of
the adjoint by the conjugate-dual polarization. In rank three
use \(F=K\) and the cusp attached above.

For the negative component at \(c\) the differential of inverse transpose contributes the minus sign. Modulo \(t^2\),
\[
 J^{-1}(1+tX)^{-t}J=1-tJ^{-1}X^tJ .
\] A
common character twist cancels in the adjoint. The
cyclotomic irreducibility proof at the start of this
section verifies the remaining hypothesis. The theorem thus
gives precisely the rational vanishing used for the type-I
host point. One can first pass to the enlarged marked frame, transport the ordinary and local genericity data
there, apply polarized adjoint vanishing, and return using restriction--corestriction in characteristic
zero. Direct application over \(K\) gives the same group.
\hfill\mbox{\cite{thorne-adjoint}}

\begin{lemma}[Adjoint Selmer vanishing for the descended systems]
\label{lem:descended-adjoint-vanishing}
With the applicable polarization and absolute irreducibility
on the cyclotomic kernel, let \(V\) be the rational adjoint of the descended system over the CM field \(F\),
extended to its maximal totally real subfield \(F^+\) by the polarization.
In the specified cases
\[
                            H_f^1(F^+,V)=0 .
\] The proof
handles the rank-two adjoint diagonal and off-diagonal
summands separately.
\end{lemma}

\begin{proof}
\proofstep{Strong irreducibility of the rational rank-two block.}
For the general-data corner it suffices to prove cyclotomic
irreducibility on \(K\). Inertia \(I\) at \(\ell\) lies
in \(H\) and has full cyclotomic image since \(K\) is
unramified at \(\ell\). Let \(H'=H\cap\ker\epsilon\). The preceding equality of cyclotomic images means that for each \(h\in H\) there is \(i\in I\) with the same image. Thus
\[
 i^{-1}h\in H',\qquad H=IH'.
\] If \(W|_{H'}\) were reducible,
the normal restriction argument would make it scalar or a
sum of two different characters. In the scalar case the
ordinary inertia line would be stable under \(I H'\).

In
the second case inertia permutes the two lines. Choose in
\(I\cap\ker\psi\) an element with infinite-order cyclotomic
value. An element that exchanged these two lines would have a matrix of the following shape in their basis:
\[
 \begin{pmatrix}0&a\\ b&0\end{pmatrix}^{\!2}=ab\,\mathrm{id}.
\]
Its eigenvalue ratio would then have square one, contrary to the infinite-order ratio of the element just chosen. Its two eigenlines are consequently the
two \(H'\)-lines. The ordinary inertia line must be one of
these, and again is stable under \(I H'\). Both alternatives
contradict irreducibility on \(H\). The rank-two cusp over
\(K\) therefore gives polarized adjoint vanishing there
by the same theorem.

\proofstep{The diagonal and the off-diagonal adjoint summands.}
Decompose the polarized endomorphisms relative to the
outer plane and the middle line. The outer endomorphisms
have this rank-two polarized action. Indeed the symmetric
pairing from the original frame on the outer plane is the
reversing matrix, equal up to scalar to the alternating
determinant pairing composed with \(W(c)^{-1}\).
On the one-dimensional middle block, transpose fixes every scalar. Thus the polarized action of \(c\) on its endomorphisms is the minus sign:
\[
 x\longmapsto-x^t=-x\qquad(x\in\operatorname{End}(\text{middle line})).
\]
This summand has character \(\eta\).

To see its finite Selmer group vanishes in characteristic zero,
observe that after restriction to \(K\) a finite class with trivial
coefficients is unramified at every finite prime (at
\(\ell\), a crystalline extension of the trivial character
by itself is unramified). The class group is finite,
and restriction is injective in characteristic zero.
The two off-diagonal spaces are interchanged by \(c\).
One restricts to \(W|_H\) and the other to \((W|_H)^\vee\).
On \(H\) the determinant pairing identifies these two modules.
Induction and the projection formula therefore give their sum as
\[
 \operatorname{Ind}_{H}^{G_\Q}(W|_H)
  \simeq W\otimes\operatorname{Ind}_{H}^{G_\Q}\mathbf 1
  \simeq W\oplus(W\otimes\eta).
\]
The second isomorphism uses that \(H\) has index two and the
coefficient field has characteristic zero.

For the newform \(f\) above put \(k=2e+1\). In Kato's
geometric Frobenius convention the representation
\(V_\lambda(f)\) is \(W(-e)\) after scalar extension.
Apply his Theorem~14.2(1) over the abelian field \(K\),
at the twist \(r=e\). Indeed \(1\le e\le k-1\) and
\(e\ne k/2\).

For every stable lattice the associated
torsion Selmer group is finite. A nonzero rational line of finite Selmer classes would give an infinite subgroup of the torsion Selmer group. Indeed for a line \(L\) over the coefficient field \(E\) and any full lattice \(\Lambda\subset L\),
\[
 L/\Lambda\simeq E/\mathcal O_E,\qquad |E/\mathcal O_E|=\infty .
\]
The torsion bound thus forces the rational finite Selmer group to vanish. Here the
integral global cohomology unramified off the finite set
of bad places and \(\ell\) is finitely generated, so its
image is a lattice. The local finite conditions are
compatible with the maps to torsion coefficients. Both
\(W\) and \(W\otimes\eta\) restrict to \(W|_H\). Finite
classes restrict to finite classes, and corestriction
composed with restriction is multiplication by two.
This proves vanishing on both off-diagonal spaces over
\(\Q\), including at \(e=1\). The direct-sum decomposition
now gives the entire polarized type-II vanishing.
\hfill\mbox{\cite{kato-modular,thorne-adjoint}}

\end{proof}

\paragraph{Fixed-inertia derivatives on the auxiliary field over \(\mathbb Q(i)\).}
The fixed-inertia condition on the Gaussian field requires a local derivative
calculation of its own. It will be used in the comparison with the split point.

\begin{lemma}[Fixed-inertia derivatives are finite]\label{lem:gaussian-fixed-inertia-finite}
For the marked Gaussian test and its prescribed ordinary
characters, let \(V\) denote the polarized adjoint and \(z\) the first-order
cocycle of a normalized dual-number deformation of the Gaussian test
representation. For every deformation satisfying the fixed-inertia equations,
\[
                       [z]\in H_f^1(\mathbb Q,V)
                              \ \subseteq\ H^1(\mathbb Q,V).
\]
\end{lemma}

\begin{proof}
A deformation to dual numbers with first derivative \(y\)
has matrix part \((1+\delta y(g))r(g)\), \(\delta^2=0\).
This defines a cocycle for the polarized adjoint. At an
away prime outside the bad set the deformation is
unramified. At a level prime outside the exceptional set the
local equations kill the derivative on inertia in
\(\ker\eta\cap\ker P\), a finite-index subgroup of inertia.

At an exceptional prime the dual-number deformation on
positive inertia is conjugate to the prescribed diagonal
action with its exact three characters. Differentiating
this equality of matrices expresses the cocycle there
as a coboundary, with a single conjugating derivative
for all inertia elements. For either finite-index inertia subgroup \(I'\subset I\), let
\(\xi\in H^1(I,V)\) be this local class over the coefficient field.
The preceding restriction calculation gives
\[
 [I:I']\,\xi
 =\operatorname{cor}_{I'}^I\operatorname{res}_{I'}^I(\xi)=0 .
\]
The positive integer \([I:I']\) is invertible over the characteristic-zero
coefficient field, so \(\xi=0\) on full inertia.

\proofstep{The finite local condition above the coefficient prime.}
At \(\ell\) restrict to the open normal subgroup of the
decomposition group killing \(\eta\) and the
\(\ell\)-local component of \(\psi\). On its inertia
the ordinary dual-number characters are exactly
\(\epsilon^{e},1,\epsilon^{-e}\). The separator
\(x^{\ell^{v_\ell(N)}}\) belongs to this subgroup since its
cyclotomic value is one modulo \(\ell^{v_\ell(N)+1}\)
by successive powering. In characteristic \(\ell\) the separator retains its distinct residual eigenvalues after powering, since
\[
 a^{\ell^k}-b^{\ell^k}=(a-b)^{\ell^k}\ne0\qquad(a\ne b).
\] Thus no two different
graded characters differ by \(\epsilon\) on this subgroup.
Viewed over the coefficient field, each dual-number line
is a self-extension of a crystalline character by an
unramified additive derivative, so is crystalline. To
add the next such two-dimensional block to the preceding
ones, the Hom space from the new block to the preceding
part has a filtration with characters of strictly
positive cyclotomic exponents.

Its dual twist has no
invariants by the separator. The local dimensions for this Hom space agree. Since the finite subspace embeds naturally in the full group,
\[
 H_f^1\subseteq H^1,\qquad \dim H_f^1=\dim H^1\quad\Longrightarrow\quad H_f^1=H^1 .
\]
The extension of crystalline blocks is therefore crystalline. This induction proves that the entire
self-extension is crystalline on the subgroup. Equivalently
the adjoint cocycle splits after tensoring with
\(B_{\rm cris}\) there. Averaging a splitting over the
finite quotient gives one over the original group. The
cocycle is therefore finite at \(\ell\) too. All matrices
and changes of basis may be taken over a finite
coefficient extension.

To see this for a continuous cocycle \(z:G\to V_{\overline{\Q}_\ell}\), fix a finite field of definition for the action. The sets \(H_L=z^{-1}(V_L)\), for the countably many finite extensions \(L\) of that field in \(\overline{\Q}_\ell\), are closed and cover \(G\). The cocycle identity makes each a subgroup. By Baire's theorem one \(H_L\) has interior, hence is open. Choose coset representatives \(g_1,\ldots,g_b\). For \(h\in H_L\),
\[
             z(g_i h)=z(g_i)+g_i z(h)\in z(g_i)+V_L .
\]
Adjoining the finitely many coordinates of the \(z(g_i)\) gives a common finite field of definition.
\hfill\mbox{\cite{bloch-kato-finite}}

\end{proof}

\begin{construction}
\[
\begin{array}{lll}
\text{branch}&\text{constituent descended}&\text{Gaussian application}\\ \hline
\text{type I}&\text{irreducible rank three}&\text{enters the trace product}\\
\text{type II}&\text{outer plane in rank two}&\text{excluded under }\mathrm{Xpp}
\end{array}
\]
\caption{The two possible branches at a two-sided point. The exclusion in the Gaussian application uses the residual vanishing assumption for the outer plane.}\label{fig:descent-trace-branches}
\end{construction}

The two branches in the table tell us which descended constituents to count. We can
now form the trace products on the original fields.

\begin{proposition}[Trace
products on the original packets]
\label{prop:original-packet-count}
For the indicated type-I and outer-plane points, the trace
packets and denominator allowances give the nonzero integer \(A_{J_0}(p)\) of \eqref{eq:local-odd-d}.
In the notation of that estimate its size satisfies
\[
 \log |A_{J_0}(p)|\ll_\psi
   \varphi([U:J_0])\bigl(\mathcal N(j+1)\bigr)^a\log p .
\] The Gaussian hypothesis excludes the separate
outer-plane case in that application. The power of
\(N(e+1)\) depends only on the original quadratic field and
fixed dimension.
\end{proposition}

\begin{proof}
\proofstep{The residual grading-torus tangent.}
The Gaussian tangent class is now zero by the type-I
vanishing. On either host a coboundary is an
infinitesimal conjugation by \(1+\delta X\). The
normalization of the marked difference forces \(X\)
diagonal because the off-diagonal entries of its commutator
with \(\operatorname{diag}(z,0,-z)\) must vanish, with
\(z\ne0\). Normalization at \(c\) imposes the adjoint condition on the surviving diagonal entries, leaving
\[
 X=-J^{-1}X^tJ,\qquad X=\operatorname{diag}(u,0,-u).
\] This is exactly the
infinitesimal grading action. At a point with nonzero
coordinates of both signs it accounts for the tangent
along its conjugacy orbit.

\proofstep{Which points enter the two products.}
A nonzero entry of each strict sign precludes both triangular
orientations. The marked coordinate test in
Section~\ref{sec:marked-hosts} then leaves type I and the
type-II outer plane. The rank-three trace estimate keeps both. For the case over the Gaussian field the
extra residual hypothesis used by the submission is vanishing
for the constrained weight \(-2\) classes, denoted
\(\mathrm{Xpp}\). In zero-based indices the coefficient on the entry \((2,0)\) is the one-dimensional module
\[
 k(\eta P^{-2}),\qquad g\cdot x=\eta(g)P(g)^{-2}x .
\] We explain exactly where this extra
arithmetic enters.

\proofstep{The residual class detected by an outer-block entry.}
In a type-II block the entry \((2,1)\) is identically zero.
If any \((2,0)\)-entry is nonzero, choose \(d=r(g_0)_{20}\)
of smallest valuation and put \(b(g)=r(g)_{20}/d\).
These quotients are integral. Multiplication on the positive
component gives, on reducing,
\[
  z(g)=\overline{b(g)}P(g)^{-1},\qquad
  z(gh)=z(g)+\eta(g)P(g)^{-2}z(h).
\]
For the extension of this equality to the negative component
use \(r(c)_{\rm mat}=1\). On the outer block the action of
\(c\) is inverse transpose with the reversing pairing, which
sends its lower left entry to its negative divided by its
determinant. That determinant is one modulo the maximal
ideal. This gives the displayed sign and \(z(c)=0\).
Continuity of the residual quotient follows by evaluating
the continuous entry modulo one power of the uniformizer
beyond \(v(d)\). Away from \(T_b\) the killed inertia
elements have zero entry.

At a place in \(T_b\), write
the first ordinary local line in the integral model with
generator \((1,h_1,h_2)^t\) congruent to the first coordinate
vector. For every positive local \(g\), let \(V(g)\) be
the eigenvalue on this vector minus \(r(g)_{22}\).
Reading the last row now gives the proportionality
\[
 d\,b(g)=h_2 V(g).
\] At the separating
element \(s\) of this place the residue of \(V(s)\) is
\(P(s)-P(s)^{-1}\ne0\). Hence \(h_2/d=b(s)/V(s)\)
is integral, and on reduction \(z(g)\) is a single
constant times \(1-P(g)^{-2}\) for all such \(g\).
This also applies at the exceptional places where the
local model is diagonal. Local coboundaries extend from
the positive kernel by averaging over its quotient of
order at most two. Thus \(z\) satisfies the indicated
residual Selmer conditions.

If the residue coefficients
were enlarged, apply each linear functional to the base
residue field before using the vanishing. Each scalar cocycle is a coboundary \(g\mapsto(\eta(g)P(g)^{-2}-1)a\). At conjugation this reads
\[
 0=z(c)=(-1-1)a=-2a .
\]
Since \(\ell\) is odd, \(a=0\) for every scalar functional, and therefore \(z=0\). This
contradicts \(b(g_0)=1\). The type-II block would therefore
itself be upper triangular. Under \(\mathrm{Xpp}\) only
type I contributes to the two-sided Gaussian points.

\proofstep{Exclusion in the Gaussian case.}
For the present auxiliary field over \(\mathbb Q(i)\) this residual vanishing follows
from the Bernoulli unit calculation in
Paragraph~\ref{par:dyadic-residual}. We give the conductor
check because the character here involves a square.
Use \(q_0\) for the bound on the conductor of the square
of the prime-to-\(\ell\) odd-conductor part of \(\psi\).
It is the fixed target conductor in the application over the Gaussian field.
Write \(\omega\) for the Teichmüller lift of the residual
cyclotomic character. Lift \(\eta P^2\) as the product of
the primitive characters inducing
\[
 \varphi=\omega^{2e}(\psi_{\rm odd}^{(\ell')})^2,
       \qquad \rho=\chi_4\psi_2^2 .
\]
Their conductors are coprime. The odd conductor \(f\) of
the first divides \(\ell q_0\). With \(d\) the order of
\(\ell\) modulo the least common multiple of four and
the order of \(\psi_{\rm odd}^{(\ell')}\), its values have
order dividing \(\ell^d-1\). Put \(r_0=v_2(\ell^d-1)\ge2\).

Since \(c=c_2(\psi)\ge4\), the value of \(\psi_2\) on
five has order \(2^{c-2}\). Squaring divides that order
by two, and \(\chi_4(5)=1\). Hence \(\rho\) has exact
conductor \(2^{c-1}\), including at \(c=4\).
Its product with \(\varphi\) is odd. The large-conductor
hypothesis gives
\[
        2^{c-1}\ge (\ell q_0)4^{r_0}\ge f4^{r_0}.
\]
The Bernoulli unit assertion applies to this product and to
its inverse. In the paragraph cited above the only
requirements for the Stickelberger argument are
unramified restriction on the kernel of the character
and the local coboundary condition at \(\ell\) to
ensure this unramifiedness there.

On
\(G_{\Q(\mu_{f2^{c-1}})}\) the inverse residual character
\(\eta P^{-2}\) is trivial. The cocycle restricts to an
unramified homomorphism because the intersection of inertia
with \(\ker\eta\cap\ker P\) has prime-to-\(\ell\) index
in inertia, and local coboundaries vanish on the character
kernel. Let \(z\) denote this restricted unramified homomorphism
and \(\Theta\) the smoothed Stickelberger element. Annihilation
and the character computation give
\[
           0=\Theta z=\chi(\Theta)z,\qquad \chi(\Theta)\in k^\times
\]
on the component being tested, with \(\chi\) its corresponding
residual character. Hence \(z=0\). The original cocycle factors
through the remaining abelian cyclotomic quotient. If \(g_0\)
has character value \(\chi_0(g_0)\ne1\) for its original
coefficient character \(\chi_0=\eta P^{-2}\), commutativity
gives
\[
   (\chi_0(g_0)-1)z(g)=(\chi_0(g)-1)z(g_0).
\]
This expresses the descended cocycle as a coboundary. This proves the required residual vanishing
on the actual host.
\hfill\mbox{\cite{washington-nonp,washington-cyclo}}

\proofstep{The trace bound for each admitted point.}
The quantitative contribution of an admitted point can
now be read at the original Frobenius. In the notation
of (D), the rank-three cusp has level \(N\) and exponents
\((e-1,-1,-e-1)\) at both complex embeddings of the fixed
\(K\). To bound its packet size polynomially in \(N(e+1)\),
use the finite cell complex at level three, the cover at level
\(3N\) and the Weyl dimension bound
in Section~\ref{sec:rank-four-deformation}. Clozel's
conjugation theorem preserves this packet and its algebraic
normalization. Let \(b\) be an algebraic Satake root at a good
place of norm \(q\). An invariant Galois lattice makes it
a unit away from the residue characteristic \(p\).
At \(p\), apply the spherical crystalline comparison and the
packet weak-admissibility estimate
(Sections~\ref{sec:local-global} and~\ref{sec:rank-four-deformation}):
\[
 v_p(q^{e-1}b^{-1})
     =v_p(q)\bigl(e-1-v_p(b)/v_p(q)\bigr)\ge0 .
\]
Indeed the maximal Hodge exponent is \(e-1\). This applies to
each coefficient embedding of each packet member.

One can also obtain the denominator bound on a
definite unitary group. After quadratic base change split
at the tested place, a character twist unramified there
makes the cusp conjugate self-dual. The Hecke operators
on the algebraic coefficient lattice bound the elementary
symmetric functions of the roots and their reciprocals.
For a monic polynomial \(X^n+c_1 X^{n-1}+\cdots+c_n\) over a nonarchimedean field, any root \(\alpha\) satisfies
\[
 |\alpha|\le\max_{1\le j\le n}|c_j|^{1/j}.
\]
Indeed a larger value would make the leading term strictly larger than every other term. Apply this bound both to the roots and to their reciprocals to obtain the slope inequalities. The shift in the weights cancels
the value of the twisting character and the split local
extension leaves \(q\) and the Satake multiset unchanged.
\hfill\mbox{\cite{clozel-motifs,borel-serre,blggt,blggt-lg}}

For the rank-two contribution one can also count the
classical newforms directly. On the rational inertia the
two characters away from \(\ell\) have conductor bounded
by \(M_2\) and the inertia image is finite. At \(\ell\)
the gap \(2e\) and the ordinary flag give potential
crystallinity by the flag argument, and the two smooth
Weil characters have the same conductor bound. A Tate twist is unramified for the Weil parameter used to compute the Artin conductor. In this normalization its exponent obeys
\[
 a(r\otimes |\cdot|^t)=a(r).
\] Thus by
local--global compatibility and the conductor formula for
newforms the conductor divides \(M_2^2\). The normalized
newforms of these conductors occur independently in
\(S_{2e+1}(\Gamma_1(M_2^2))\).

At modulus \(M_2^2\) the modular curve index is bounded by the square of that modulus:
\[
 [\mathrm{SL}_2(\mathbb Z):\Gamma_1(M_2^2)]\le (M_2^2)^2=M_2^4 .
\]
The valence dimension bound is linear in the weight times
this index, so the dimension is \(O((e+1)M_2^4)\). The
list is stable under conjugation. At a good \(p\) let
\(\alpha,\beta\) be the algebraic Hecke roots before the
twist. They are integral, have absolute value \(p^e\)
in every embedding and satisfy \(\alpha\beta=p^{2e}\xi\)
with \(\xi\) a root of unity.

If \(a=\alpha/p^e\) and
\(b=\beta/p^e\), then
\[
             p^e a=\alpha,\qquad p^e a^{-1}=\xi^{-1}\beta,
 \qquad      p^e b=\beta,\qquad p^e b^{-1}=\xi^{-1}\alpha .
\]
At residue degree \(h\), each of \(a^{\pm h},b^{\pm h}\)
is consequently integral after multiplication by \(p^{eh}\).
This factor also preserves the middle eigenvalue \(1\), and
clears the denominator in the trace.
The packet can also be counted from its irreducible restrictions to the quadratic field, the choices of roots at positive residue degree and the at most two descent characters for each fixed irreducible restriction.
The direct count gives a polynomial bound of the same
required form. \hfill\mbox{\cite{deligne-galois,carayol-local,deligne-weil}}

\proofstep{Counting isomorphism classes by torus orbits.}
For a fixed character and either type, equality of complete
systems gives equality of semisimple representations at
all good Frobenii and hence isomorphism over \(K\).
Assign a packet index to each selected conjugacy class. The complete system, rather than one trace, gives the implication
\[
 \text{same packet index}\quad\Longrightarrow\quad\text{same selected conjugacy class},
\]
proving injectivity for the count. Thus no
identification of such indices is needed. A common integer product over
both packets and the full character orbit is the integer
in (D). For the Gaussian case use just the first packet
and allow its two field embeddings, as in the case over the Gaussian field
product. In (D) the implied constant and exponent depend
only on the chosen quadratic field and the fixed dimension.
In particular neither the marked field, coefficient
ramification, omitted primes, nor the choice of ordinary
seed enters the power of \(N(e+1)\).

Section~\ref{sec:rank-three-length} supplies the integral
chart comparison for (L). The trace is a function of
weight zero. The length comparison uses one point
from each two-sided torus orbit, and the marked
bases make those representatives give distinct
restriction isomorphism classes in the original packets.
\end{proof}
